% Options for packages loaded elsewhere
\PassOptionsToPackage{unicode}{hyperref}
\PassOptionsToPackage{hyphens}{url}
\documentclass[
  indonesian,
  11pt,
  a4paper,
  oneside,
  openany]{report}
\usepackage{xcolor}
\usepackage{amsmath,amssymb}
\setcounter{secnumdepth}{-\maxdimen} % remove section numbering
\usepackage{iftex}
\ifPDFTeX
  \usepackage[T1]{fontenc}
  \usepackage[utf8]{inputenc}
  \usepackage{textcomp} % provide euro and other symbols
\else % if luatex or xetex
  \usepackage{unicode-math} % this also loads fontspec
  \defaultfontfeatures{Scale=MatchLowercase}
  \defaultfontfeatures[\rmfamily]{Ligatures=TeX,Scale=1}
\fi
\usepackage{lmodern}
\ifPDFTeX\else
  % xetex/luatex font selection
  \setmainfont[]{Charter}
  \setsansfont[]{Arial}
  \setmonofont[]{Menlo}
\fi
% Use upquote if available, for straight quotes in verbatim environments
\IfFileExists{upquote.sty}{\usepackage{upquote}}{}
\IfFileExists{microtype.sty}{% use microtype if available
  \usepackage[]{microtype}
  \UseMicrotypeSet[protrusion]{basicmath} % disable protrusion for tt fonts
}{}
\makeatletter
\@ifundefined{KOMAClassName}{% if non-KOMA class
  \IfFileExists{parskip.sty}{%
    \usepackage{parskip}
  }{% else
    \setlength{\parindent}{0pt}
    \setlength{\parskip}{6pt plus 2pt minus 1pt}}
}{% if KOMA class
  \KOMAoptions{parskip=half}}
\makeatother
\usepackage{longtable,booktabs,array}
\usepackage{caption}
\captionsetup[table]{skip=6pt}
\newcounter{none} % for unnumbered tables
\usepackage{calc} % for calculating minipage widths
% Correct order of tables after \paragraph or \subparagraph
\usepackage{etoolbox}
\makeatletter
\patchcmd\longtable{\par}{\if@noskipsec\mbox{}\fi\par}{}{}
\makeatother
% Allow footnotes in longtable head/foot
\IfFileExists{footnotehyper.sty}{\usepackage{footnotehyper}}{\usepackage{footnote}}
\makesavenoteenv{longtable}
\ifLuaTeX
\usepackage[bidi=basic,shorthands=off,provide=*]{babel}
\else
\usepackage[bidi=default,shorthands=off,provide=*]{babel}
\fi
\ifPDFTeX
\else
\babelfont{rm}[]{Charter}
\fi
\ifLuaTeX
  \usepackage{selnolig} % disable illegal ligatures
\fi
\setlength{\emergencystretch}{3em} % prevent overfull lines
\providecommand{\tightlist}{%
  \setlength{\itemsep}{0pt}\setlength{\parskip}{0pt}}
% Gaya formal monokrom untuk draf Bab 3.
\usepackage[a4paper,top=25mm,bottom=25mm,left=30mm,right=25mm,headheight=16pt,headsep=7mm,footskip=12mm]{geometry}
\usepackage{microtype}
\usepackage{setspace}
\usepackage{ragged2e}
\usepackage{titlesec}
\usepackage{enumitem}
\usepackage{needspace}
\usepackage{etoolbox}
\usepackage{array}
\usepackage{booktabs}
\usepackage{longtable}
\usepackage{tabularx}
\usepackage{adjustbox}
\usepackage{float}
\usepackage{fancyhdr}
\usepackage{newunicodechar}
\newunicodechar{∅}{\ensuremath{\varnothing}}

\definecolor{Ink}{HTML}{222222}
\definecolor{RuleGray}{HTML}{666666}
\definecolor{LinkGray}{HTML}{222222}

\AtBeginDocument{%
  \hypersetup{
    linkcolor=Ink,
    citecolor=Ink,
    urlcolor=LinkGray,
    pdfborder={0 0 0}
  }%
  \urlstyle{same}%
}

\setstretch{1.18}
\setlength{\parindent}{1.1em}
\setlength{\parskip}{0.28em}
\setlength{\emergencystretch}{2.5em}
\raggedbottom
\clubpenalty=10000
\widowpenalty=10000
\displaywidowpenalty=10000

\setlist[itemize]{leftmargin=1.7em,itemsep=0.25em,topsep=0.35em}
\setlist[enumerate]{leftmargin=1.9em,itemsep=0.25em,topsep=0.35em}

\titleformat{\chapter}[display]
  {\normalfont\bfseries\centering\color{Ink}}
  {}
  {0pt}
  {\Large\hyphenpenalty=10000\exhyphenpenalty=10000}
\titlespacing*{\chapter}{0pt}{-6pt}{18pt}
\titleformat{\section}{\Needspace{4\baselineskip}\large\bfseries\color{Ink}}{\thesection}{0.7em}{}
\titleformat{\subsection}{\Needspace{3\baselineskip}\normalsize\bfseries\color{Ink}}{\thesubsection}{0.7em}{}
\titleformat{\subsubsection}{\Needspace{3\baselineskip}\normalsize\bfseries\color{Ink}}{\thesubsubsection}{0.7em}{}

\pagestyle{plain}
\fancyhf{}
\fancyfoot[C]{\thepage}
\renewcommand{\headrulewidth}{0pt}
\renewcommand{\footrulewidth}{0pt}
\fancypagestyle{plain}{%
  \fancyhf{}
  \fancyfoot[C]{\thepage}
  \renewcommand{\headrulewidth}{0pt}
}

\renewcommand{\arraystretch}{1.28}
\setlength{\extrarowheight}{1pt}
\setlength{\tabcolsep}{4pt}
\setlength{\LTpre}{0.5em}
\setlength{\LTpost}{0.8em}
\AtBeginEnvironment{longtable}{\footnotesize\setlength{\tabcolsep}{3pt}\setlength{\extrarowheight}{1pt}}

\AtBeginEnvironment{quote}{\small\setlength{\leftskip}{1.2em}\setlength{\rightskip}{1.2em}}
\usepackage{bookmark}
\IfFileExists{xurl.sty}{\usepackage{xurl}}{} % add URL line breaks if available
\urlstyle{same}
% fallback for those not using the hyperref driver hyperxmp:
\makeatletter
\@ifundefined{xmpquote}{\newcommand{\xmpquote}[1]{#1}}{}
\makeatother
\hypersetup{
  pdftitle={Bab 3 - Metode Penelitian},
  pdflang={id-ID},
  hidelinks,
  pdfcreator={LaTeX via pandoc}}

\title{Bab 3 - Metode Penelitian}
\usepackage{etoolbox}
\makeatletter
\providecommand{\subtitle}[1]{% add subtitle to \maketitle
  \apptocmd{\@title}{\par {\large #1 \par}}{}{}
}
\makeatother
\subtitle{Draf rancangan metode Borrowed Size dan Agglomeration Shadow
pada Kawasan Aglomerasi Jakarta}
\author{}
\date{26 Agustus 2026}

\begin{document}
\maketitle

\renewcommand*\contentsname{Daftar Isi}
{
\setcounter{tocdepth}{1}
\tableofcontents
}
\chapter{BAB 3 METODE PENELITIAN}\label{bab-3-metode-penelitian}

\section{3.1 Pendekatan penelitian}\label{pendekatan-penelitian}

\subsection{3.1.1 Pilihan pendekatan}\label{pilihan-pendekatan}

Penelitian ini menggunakan \textbf{pendekatan kuantitatif (deduktif)
dengan analisis spasial-relasional}. Konstruk integrasi metropolitan,
fungsi relatif, manfaat, beban, serta gerbang penggunaan label
\emph{borrowed size} dan \emph{agglomeration shadow} diturunkan dari
kerangka konseptual Bab 2 sebelum klasifikasi dilakukan. Pendekatan ini
dipilih karena pertanyaan penelitian menuntut pengukuran hubungan antara
pusat sekunder dan Jakarta, perbandingan fungsi aktual dengan fungsi
yang diharapkan dari ukuran lokal, serta pengujian stabilitas hasil pada
variasi asumsi.

Pendekatan ini tidak berarti semua konstruk harus direduksi menjadi satu
angka. Integrasi, fungsi relatif, manfaat, ketergantungan, dan beban
dipertahankan sebagai dimensi yang dapat dibaca bersama. Statistik
deskriptif, pemetaan, analisis jaringan, tolok ukur, residual, dan uji
sensitivitas digunakan sesuai jenis bukti yang tersedia.

Audit dokumen dan audit asal-usul data merupakan prosedur penjaminan
mutu. Prosedur ini digunakan untuk memeriksa definisi, tahun, unit,
cakupan, cara pembentukan, nilai hilang, dan batas penggunaan data.
Pemeriksaan dokumen tersebut tidak menjadikan penelitian ini penelitian
campuran karena dokumen tidak dianalisis sebagai data untuk menjawab
pertanyaan substantif dan tidak digunakan untuk mengisi nilai numerik
yang tidak tersedia.

\subsection{3.1.2 Alasan kesesuaian dengan pertanyaan
penelitian}\label{alasan-kesesuaian-dengan-pertanyaan-penelitian}

Kesesuaian pendekatan dengan pertanyaan penelitian diringkas sebagai
berikut.

{\def\LTcaptype{none} % do not increment counter
\begin{longtable}[]{@{}
  >{\raggedright\arraybackslash}p{(\linewidth - 4\tabcolsep) * \real{0.3333}}
  >{\raggedright\arraybackslash}p{(\linewidth - 4\tabcolsep) * \real{0.3333}}
  >{\raggedright\arraybackslash}p{(\linewidth - 4\tabcolsep) * \real{0.3333}}@{}}
\toprule\noalign{}
\begin{minipage}[b]{\linewidth}\raggedright
Kode
\end{minipage} & \begin{minipage}[b]{\linewidth}\raggedright
Fokus pertanyaan
\end{minipage} & \begin{minipage}[b]{\linewidth}\raggedright
Alasan menggunakan pendekatan kuantitatif-spasial
\end{minipage} \\
\midrule\noalign{}
\endhead
\bottomrule\noalign{}
\endlastfoot
Q1 & Derajat, arah, dan beban hubungan pusat sekunder dengan Jakarta
dan, jika bukti mendukung, dengan pusat lain & Hubungan perlu diukur
pada unit asal-tujuan atau unit orientasi yang sah dan dibandingkan
secara spasial. \\
Q2 & Fungsi atau kinerja lokal dibandingkan dengan tingkat yang
diharapkan berdasarkan ukuran serta karakteristik lokal & Perbandingan
memerlukan indikator fungsi dan tolok ukur yang menghasilkan surplus
atau defisit relatif per konstruk. \\
Q3 & Kombinasi integrasi, fungsi relatif, manfaat, dan beban pada
\emph{status quo} & Dimensi perlu dibaca sebagai profil dan gradien
sebelum diringkas menjadi kategori. \\
Q4 & Ketahanan diagnosis terhadap indikator, unit, bobot, ambang, dan
status mutu bukti & Hasil perlu dihitung ulang dengan spesifikasi
terdokumentasi tanpa mengubah nilai observasi garis dasar. \\
\end{longtable}
}

Pendekatan ini tidak cukup untuk menyatakan dampak kausal pembangunan
jalan, perluasan wilayah, atau kebijakan tertentu tanpa strategi
identifikasi tambahan. Karena itu, istilah yang dipakai dalam Bab 3
adalah \emph{asosiasi}, \emph{pola}, \emph{surplus atau defisit
relatif}, \emph{ketergantungan yang konsisten dengan data}, dan
\emph{ketahanan klasifikasi}, bukan ``dampak kausal'' kecuali desain
kelak benar-benar diperkuat.

\subsection{3.1.3 Hubungan dengan skenario
metodologis}\label{hubungan-dengan-skenario-metodologis}

Skenario akses data tidak mengganti pertanyaan substantif. Skenario
menentukan kedalaman pengukuran dan kekuatan klaim. Data yang lebih
rinci dapat mengaktifkan metode yang lebih kuat; data yang lebih kasar
mengharuskan unit, indikator, atau klaim diturunkan. Dengan demikian,
penelitian tidak memilih tesis baru setiap kali satu pintu data gagal.

\section{3.2 Desain penelitian}\label{desain-penelitian}

\subsection{3.2.1 Bentuk desain}\label{bentuk-desain}

Desain penelitian adalah \textbf{studi kasus tunggal kuantitatif dengan
potret potong lintang antarruang, diagnosis relasional, dan analisis
sensitivitas}. Kasus tunggal dipilih agar pusat-pusat sekunder dapat
dibandingkan di dalam satu sistem metropolitan tanpa memperlakukan
perbedaan antarmetropolitan sebagai variasi yang setara. Sabuk
pembanding fungsional hanya digunakan untuk memperluas tolok ukur
apabila definisi, unit, dan periodenya kompatibel.

Desain memiliki tiga lapisan:

\begin{enumerate}
\def\labelenumi{\arabic{enumi}.}
\tightlist
\item
  \textbf{Lapisan observasi:} data langsung, agregat, sintetis, atau
  proksi pada unit dan periode yang tercantum dalam manifest.
\item
  \textbf{Lapisan diagnosis:} pengukuran integrasi, fungsi relatif,
  manfaat, beban, dan tipologi dari satu konfigurasi bukti yang
  dibekukan.
\item
  \textbf{Lapisan ketahanan dan eksplorasi:} perubahan konfigurasi
  bukti, spesifikasi model, atau nilai skenario yang dibatasi dan selalu
  dilaporkan terhadap garis dasar tanpa menimpa data observasi.
\end{enumerate}

\subsection{3.2.2 Urutan kerja
penelitian}\label{urutan-kerja-penelitian}

Urutan penelitian sementara adalah sebagai berikut.

\begin{enumerate}
\def\labelenumi{\arabic{enumi}.}
\tightlist
\item
  Menetapkan pertanyaan, konstruk, definisi kerja, dan aturan klaim
  sebelum melihat hasil klasifikasi.
\item
  Menginventarisasi berkas dan mencatat asal-usul data setiap modul.
\item
  Mengaudit tahun, unit, cakupan, definisi, nilai hilang, reliabilitas,
  dan batas publikasi.
\item
  Mengharmonisasi geometri dan tabel korespondensi tanpa melakukan
  disagregasi semu.
\item
  Menetapkan kandidat pusat sekunder dan wilayah tangkapan fungsional
  berdasarkan bukti fungsi atau hubungan yang tersedia.
\item
  Mengukur integrasi metropolitan sesuai tingkat bukti M0-M3.
\item
  Mengukur fungsi lokal dan fungsi relatif sesuai tingkat bukti U0-U3.
\item
  Memisahkan profil manfaat atau \emph{functional upgrading} dari profil
  beban atau ketergantungan.
\item
  Menyusun diagnosis \emph{status quo} sebagai gradien dan, bila
  diperlukan, ringkasan kelas.
\item
  Menjalankan diagnostik spasial dan uji sensitivitas.
\item
  Menerjemahkan model yang sama ke WebGL dengan garis dasar, skenario,
  selisih, dan peta ketahanan.
\item
  Mengunci manifest, tabel indikator, parameter, dan batas klaim untuk
  reproduksi.
\end{enumerate}

\subsection{3.2.3 Kelebihan dan keterbatasan
desain}\label{kelebihan-dan-keterbatasan-desain}

Kelebihan desain ini adalah transparan terhadap variasi data, dapat
berjalan dengan garis dasar terbuka, dan tidak memaksakan satu label
biner. Desain juga memungkinkan data permohonan dan data terbuka
dibandingkan sebagai dua konfigurasi bukti.

Keterbatasannya adalah potret potong lintang tidak membuktikan perubahan
temporal, proksi tidak sama dengan konstruk langsung, dan unit analisis
yang berbeda dapat mengubah hasil. Keterbatasan tersebut dikelola
melalui pelabelan status bukti, analisis multiresolusi, tolok ukur yang
eksplisit, uji sensitivitas, dan pembatasan klaim; prosedur tersebut
tidak menghapus keterbatasan dasar desain.

\subsection{3.2.4 Perbedaan diagnosis dan
skenario}\label{perbedaan-diagnosis-dan-skenario}

Nilai \emph{status quo} dibekukan sebagai garis dasar. Perubahan
penggeser atau parameter tidak boleh menimpa nilai observasi, mengubah
batas administratif, atau disajikan sebagai prediksi. Setiap skenario
menyimpan konfigurasi indikator, bobot, ambang, nilai yang diubah, dan
alasan substantif perubahan.

Perubahan karena sumber data diperoleh, hilang, atau turun mutu dicatat
sebagai \textbf{skenario bukti}. Perubahan bobot, ambang, atau
spesifikasi dalam konfigurasi bukti yang sama dicatat sebagai
\textbf{skenario parameter}. Perubahan nilai indikator untuk pertanyaan
\emph{what-if} merupakan eksplorasi kondisional tambahan. Ketiganya
tidak boleh ditampilkan sebagai jenis perubahan yang sama.

\section{3.3 Ruang lingkup penelitian}\label{ruang-lingkup-penelitian}

\subsection{3.3.1 Ruang lingkup spasial}\label{ruang-lingkup-spasial}

Kasus utama adalah Jabodetabek sebagai kawasan utama penelitian, dengan
sabuk Wilayah Statistik Metropolitan sebagai pembanding fungsional
apabila datanya kompatibel. Kasus ini dipilih karena memungkinkan
hubungan antara satu inti dan kandidat pusat sekunder lintas batas
administratif diperiksa di dalam sistem metropolitan yang sama. DKI
Jakarta diperlakukan sebagai satu inti untuk hubungan eksternal, dengan
batas administratif tetap. Data rinci di dalam DKI tetap dapat digunakan
untuk pemeriksaan internal, tetapi tidak mengubah keputusan bahwa DKI
diperlakukan sebagai satu inti dalam klasifikasi eksternal.

Wilayah di luar DKI tidak disebut sebagai ``wilayah Jakarta'' hanya
karena masuk dalam gradien analisis. Istilah yang digunakan adalah
\textbf{jejak fungsional Jakarta}, \textbf{zona analitis keterkaitan
metropolitan}, atau \textbf{gradien integrasi-ketergantungan
metropolitan}. Istilah tersebut bukan zonasi legal RTRW dan tidak
mengubah kewenangan pemerintah daerah.

Seluruh wilayah administratif di sekitar Jakarta tidak diasumsikan
otomatis menjadi unit analisis yang tepat. Delineasi akhir harus
menjelaskan sumber, tahun, definisi, dan alasan memasukkan atau
mengeluarkan setiap unit.

\subsection{3.3.2 Ruang lingkup temporal}\label{ruang-lingkup-temporal}

Garis dasar dominan direncanakan pada 2024. Data 2023, 2025, atau tahun
lain dipertahankan sebagai observasi berlabel dan tidak dilebur
seolah-olah berasal dari tahun yang sama. Jika data beberapa tahun
memiliki definisi dan unit yang sebanding, data tersebut dapat dipakai
untuk membaca perubahan atau validasi temporal.

Potret satu tahun hanya menunjukkan posisi relatif pada periode
pengamatan. Ia tidak otomatis membuktikan \emph{functional upgrading},
kematangan dari waktu ke waktu, atau titik balik. Klaim temporal hanya
digunakan jika tersedia rangkaian data dan aturan harmonisasi yang
memadai.

\subsection{3.3.3 Ruang lingkup
substantif}\label{ruang-lingkup-substantif}

Penelitian mencakup:

\begin{itemize}
\tightlist
\item
  integrasi atau eksposur metropolitan;
\item
  fungsi lokal dan fungsi relatif terhadap ukuran serta karakteristik
  lokal;
\item
  manfaat akses, \emph{functional upgrading}, dan potensi retensi
  manfaat;
\item
  beban komuter, ketergantungan, kebocoran nilai, dan defisit fungsi;
\item
  tipologi \emph{borrowed size}, campuran/transisional,
  \emph{agglomeration shadow}, relatif independen, dan periferi lemah;
\item
  sensitivitas indikator, bobot, ambang, unit, dan status mutu bukti;
  serta
\item
  penyajian hasil dalam atlas, tabel, peta, dan WebGL.
\end{itemize}

Penelitian tidak mencakup penetapan radius geometris Jakarta,
rekomendasi perubahan batas administrasi, pembuktian kausal seluruh
kebijakan metropolitan, atau klaim produktivitas penuh jika fungsi
ekonomi pada unit yang sesuai tidak tersedia.

\section{3.4 Unit amatan dan unit
analisis}\label{unit-amatan-dan-unit-analisis}

\subsection{3.4.1 Unit amatan}\label{unit-amatan}

Unit amatan adalah satuan tempat setiap observasi dicatat pada sumber.
Satuan tersebut dapat berupa desa/kelurahan, kecamatan, kabupaten/kota,
zona asal-tujuan, titik fasilitas, atau sel raster. Kecamatan menjadi
sasaran unit amatan utama untuk modul yang memiliki bukti relasional dan
fungsi yang kompatibel karena cukup rinci untuk membedakan kandidat
pusat sekunder dan wilayah sekitarnya.

Kecamatan tidak dipaksakan ketika data kritis hanya tersedia pada
kabupaten/kota atau zona model. Dalam keadaan tersebut, perhitungan
dilakukan pada unit sumber atau sebagai analisis paralel multiresolusi.
Nilai kabupaten/kota tidak diturunkan secara artifisial menjadi variasi
kecamatan, kisi, atau titik.

Kisi atau heksagon hanya digunakan untuk data yang berasal dari raster
atau titik, misalnya cahaya malam, morfologi terbangun, atau titik
minat. Kisi tersebut merupakan lapisan pendukung pada resolusi aslinya,
bukan unit pengganti untuk data administratif yang lebih kasar.

\subsection{3.4.2 Unit analisis
substantif}\label{unit-analisis-substantif}

Unit analisis substantif adalah kandidat pusat sekunder beserta wilayah
tangkapan fungsionalnya. Secara operasional, satu pusat dapat
direpresentasikan oleh satu kecamatan, beberapa unit yang bersebelahan,
atau zona sumber selama aturan pembentukannya dibekukan dan dapat
direproduksi. Kabupaten/kota tetap menjadi lapisan agregasi kebijakan,
layanan, APBD, dan tanggung jawab pemerintahan, tetapi tidak otomatis
menjadi pusat sekunder.

Daftar kandidat awal mengikuti pusat yang telah dikenali dalam dokumen
perencanaan atau sumber resmi, kemudian diuji dengan bukti kepadatan
aktivitas, fungsi layanan, pekerjaan, kawasan industri, posisi jaringan,
atau arus yang tersedia. Variabel untuk delineasi dipisahkan dari
variabel hasil sejauh data memungkinkan. Jika indikator yang sama harus
digunakan pada kedua tahap, ketahanan klasifikasi diuji kembali dengan
mengeluarkan indikator tersebut agar penetapan pusat tidak menjadi
pembuktian sirkular.

\subsection{3.4.3 Aturan keterbandingan
unit}\label{aturan-keterbandingan-unit}

Setiap indikator dicatat dengan unit asli, unit perhitungan, unit
pelaporan, dan operasi transformasinya. Satu klasifikasi hanya dihitung
jika bukti relasional dan bukti fungsi dapat dipadankan pada unit,
wilayah tangkapan, dan periode yang sebanding atau dapat diagregasikan
secara sah. Harmonisasi mengikuti aturan:

\begin{enumerate}
\def\labelenumi{\arabic{enumi}.}
\tightlist
\item
  gunakan unit asli untuk perhitungan awal;
\item
  gunakan tabel korespondensi resmi jika tersedia;
\item
  agregasikan unit rinci ke unit lebih kasar bila operasi tersebut sah;
\item
  jalankan analisis paralel bila unit tidak dapat disetarakan;
\item
  keluarkan modul jika keterbandingan tidak dapat dipertanggungjawabkan;
  dan
\item
  jangan membandingkan skor lintas konfigurasi bukti sebagai perubahan
  kondisi tanpa menyatakan perbedaan sumber dan unitnya.
\end{enumerate}

\subsection{3.4.4 Populasi, cakupan, dan validasi
manual}\label{populasi-cakupan-dan-validasi-manual}

Populasi analisis adalah seluruh kandidat pusat dan wilayah tangkapan
yang masuk dalam delineasi akhir serta mempunyai bukti yang memenuhi
aturan modul. Jika sumber merupakan sensus atau inventaris seluruh unit,
penelitian tidak mengambil sampel statistik tambahan. Pemeriksaan manual
dilakukan secara purposif pada tiga sampai lima pusat sekunder prioritas
untuk memeriksa lokasi kampus dan rumah sakit; titik minat lain
diperiksa melalui audit silang antardataset, bukan satu per satu.
Kriteria pemilihan pusat validasi dicatat sebelum pemeriksaan. Validasi
terbatas tersebut bukan sampel inferensial dan tidak digeneralisasikan
sebagai sensus seluruh kawasan.

\section{3.5 Metode pengumpulan data}\label{metode-pengumpulan-data}

\subsection{3.5.1 Prinsip pengumpulan
data}\label{prinsip-pengumpulan-data}

Pengumpulan data berjalan melalui dua jalur paralel:

\begin{enumerate}
\def\labelenumi{\arabic{enumi}.}
\tightlist
\item
  \textbf{Garis dasar terbuka}, agar prototipe dan analisis minimum
  tidak tertahan oleh permohonan data.
\item
  \textbf{Permohonan atau akses kelembagaan}, untuk mengganti,
  mengkalibrasi, atau memvalidasi proksi ketika keluaran yang diterima
  kompatibel.
\end{enumerate}

Persetujuan permohonan tidak sama dengan kelayakan analisis. Isi berkas,
definisi, unit, tahun, cakupan, dan mutu bukti menentukan status
penggunaannya.

\subsection{3.5.2 Kelompok sumber data}\label{kelompok-sumber-data}

Tabel 3.1 memuat kelompok sumber yang direncanakan, bukan pernyataan
bahwa seluruh berkas telah diterima. Sumber terbuka, sumber yang
memerlukan akses kelembagaan, dan sumber yang baru dapat dinilai setelah
berkas diterima tetap dibedakan dalam registri akses.

{\def\LTcaptype{none} % do not increment counter
\begin{longtable}[]{@{}
  >{\raggedright\arraybackslash}p{(\linewidth - 6\tabcolsep) * \real{0.2500}}
  >{\raggedright\arraybackslash}p{(\linewidth - 6\tabcolsep) * \real{0.2500}}
  >{\raggedright\arraybackslash}p{(\linewidth - 6\tabcolsep) * \real{0.2500}}
  >{\raggedright\arraybackslash}p{(\linewidth - 6\tabcolsep) * \real{0.2500}}@{}}
\toprule\noalign{}
\begin{minipage}[b]{\linewidth}\raggedright
Modul
\end{minipage} & \begin{minipage}[b]{\linewidth}\raggedright
Sumber yang direncanakan
\end{minipage} & \begin{minipage}[b]{\linewidth}\raggedright
Peran metodologis
\end{minipage} & \begin{minipage}[b]{\linewidth}\raggedright
Batas awal
\end{minipage} \\
\midrule\noalign{}
\endhead
\bottomrule\noalign{}
\endlastfoot
Batas dan unit & BIG, API geospasial BPS, administrasi resmi & geometri,
kode, tabel korespondensi, batas DKI tetap & versi dan tahun harus
dicatat \\
Ukuran lokal & WSM 2024, SP2020, publikasi statistik daerah & penduduk,
luas, kepadatan, karakteristik lokal & sumber penduduk dapat berbeda \\
Mobilitas & Survei Komuter Jabodetabek 2023, WSM 2024 & asal-tujuan
(OD), orientasi, durasi, biaya, beban & resolusi dan bobot menentukan
klaim \\
Pekerjaan & Sakernas Agustus 2024, tabulasi instansi & pekerjaan menurut
lokasi dan sektor & unit kecamatan belum pasti \\
Fungsi lokal & Podes 2024, fasilitas resmi, titik minat (POI), direktori
& layanan, fasilitas, kehadiran fungsi & POI bukan sensus fungsi \\
Transportasi & BPTJ/DJITM, JUTPI 2/3, GTFS, OSM & arus,
bangkitan-tarikan, jaringan, aksesibilitas & model atau marginal bukan
OD aktual \\
Ekonomi dan industri & PDRB, investasi, kawasan industri, direktori
usaha & konteks ekonomi dan spesialisasi & agregat tidak diturunkan ke
kecamatan \\
Lingkungan aktivitas & morfologi terbangun, VIIRS, raster terbuka &
validasi silang intensitas aktivitas & bukan ukuran produktivitas
langsung \\
\end{longtable}
}

Enam pintu permohonan utama adalah Podes 2024, Survei Komuter
Jabodetabek 2023, Sakernas Agustus 2024, Kemenhub/DJITM, JUTPI 2/3, dan
Disnaker/DPMPTSP. Jawaban Disnaker dan DPMPTSP dicatat terpisah setelah
diterima karena cakupan, definisi, unit, dan kelengkapannya dapat
berbeda.

Klasifikasi inti-periferi WSM 2024 belum diperlakukan sebagai kebenaran
dasar karena dokumen yang telah diaudit memuat ketidakkonsistenan ambang
wilayah penyangga dan rujukan tahun penduduk. Sampai klarifikasi
diperoleh, persentase orientasi yang tersedia dapat dinilai sebagai
bukti M2 pada unit sumber, sedangkan status inti-periferi hanya
digunakan setelah versi, ambang, dan tahun rujukannya dijelaskan. Jika
tabel yang diperlukan tersedia, kedua ambang yang tercantum diuji
sebagai sensitivitas delineasi.

\subsection{3.5.3 Instrumen pengumpulan dan dokumentasi
data}\label{instrumen-pengumpulan-dan-dokumentasi-data}

Karena penelitian menggunakan data sekunder, instrumen pengumpulan bukan
kuesioner atau panduan wawancara. Instrumen kerja terdiri atas:

\begin{enumerate}
\def\labelenumi{\arabic{enumi}.}
\tightlist
\item
  log akses untuk mencatat alamat sumber, tanggal akses, dan syarat
  penggunaan;
\item
  formulir audit berkas dan registri asal-usul data;
\item
  kamus data yang memuat definisi variabel, unit, tahun, nilai hilang,
  dan polaritas;
\item
  tabel korespondensi kode serta geometri antarversi;
\item
  manifest konfigurasi bukti dan parameter; serta
\item
  lembar pemeriksaan manual untuk kampus dan rumah sakit pada pusat
  prioritas.
\end{enumerate}

Skrip akuisisi, pembersihan, transformasi, dan analisis diperlakukan
sebagai bagian dari instrumen yang dapat direproduksi. Nama perangkat
lunak, versi, parameter, dan tanggal pengolahan dicatat setelah
perangkat benar-benar digunakan; Bab 3 tidak mengunci perangkat yang
belum dipakai.

\subsection{3.5.4 Audit dan status mutu
bukti}\label{audit-dan-status-mutu-bukti}

Setiap berkas dicatat dalam registri asal-usul data sekurang-kurangnya
menurut sumber, nama/versi berkas, tanggal akses, tahun observasi, unit,
cakupan, definisi, metode pembentukan, nilai hilang, reliabilitas,
lisensi, dan pembatasan publikasi.

{\def\LTcaptype{none} % do not increment counter
\begin{longtable}[]{@{}
  >{\raggedright\arraybackslash}p{(\linewidth - 4\tabcolsep) * \real{0.3333}}
  >{\raggedright\arraybackslash}p{(\linewidth - 4\tabcolsep) * \real{0.3333}}
  >{\raggedright\arraybackslash}p{(\linewidth - 4\tabcolsep) * \real{0.3333}}@{}}
\toprule\noalign{}
\begin{minipage}[b]{\linewidth}\raggedright
Status
\end{minipage} & \begin{minipage}[b]{\linewidth}\raggedright
Nama
\end{minipage} & \begin{minipage}[b]{\linewidth}\raggedright
Aturan penggunaan
\end{minipage} \\
\midrule\noalign{}
\endhead
\bottomrule\noalign{}
\endlastfoot
D & Langsung & Konstruk teramati pada unit dan periode sasaran dengan
mutu memadai. \\
A & Agregat & Konstruk teramati, tetapi unit, waktu, kategori, atau
cakupan lebih kasar. \\
S & Sintetis & Nilai diestimasi dari data lain; asumsi, kalibrasi, dan
ketidakpastian wajib ditampilkan. \\
P & Proksi & Mengukur gejala terkait, bukan konstruk yang hendak diklaim
secara langsung. \\
∅ & Dikeluarkan & Tidak tersedia, tidak sebanding, tidak reliabel, atau
tidak dapat dipertanggungjawabkan. \\
\end{longtable}
}

Status D pada tingkat permohonan tidak otomatis berarti D pada unit
penelitian. Misalnya, survei kabupaten/kota dapat berstatus A untuk
penelitian berkecamatan; jaringan dapat berstatus P untuk integrasi
aktual; dan tabel bangkitan-tarikan dapat berstatus S atau A jika tidak
memuat pasangan arus teramati.

\subsection{3.5.5 Harmonisasi spasial dan
temporal}\label{harmonisasi-spasial-dan-temporal}

Harmonisasi dilakukan setelah audit, bukan sebelum audit. Tahun 2024
dipakai sebagai jangkar utama bila indikator dapat dibandingkan.
Variabel dari 2023 atau 2025 tidak dinormalisasi ulang untuk
menyembunyikan perbedaan tahun. Setiap peta, tabel, dan skor menyimpan
tahun sumber dan status perbandingannya.

Untuk data spasial, geometri unit asli dipertahankan dalam arsip kerja.
Tabel korespondensi resmi digunakan bila tersedia. Agregasi dari desa ke
kecamatan dapat dilakukan bila definisi, kode, dan cakupan mendukung.
Disagregasi nilai kabupaten/kota ke kecamatan dilarang kecuali ada model
eksplisit, asumsi terlihat, validasi, dan pelabelan sintetis.

\subsection{3.5.6 Jalur garis dasar dan jalur
penguatan}\label{jalur-garis-dasar-dan-jalur-penguatan}

Garis dasar terbuka sekurang-kurangnya dapat memuat WSM 2024, tabel
marginal bangkitan-tarikan BPTJ, jaringan dan GTFS yang tersedia, batas
administrasi, penduduk, PDRB kabupaten/kota sebagai konteks, fasilitas
resmi, serta proksi OSM/Overture, morfologi terbangun, dan sumber
terbuka lain yang lolos audit.

Jika data OD rinci, pekerjaan menurut lokasi kerja, atau fungsi lokal
yang kompatibel diperoleh, modul tersebut digunakan untuk memperkuat
garis dasar. Data permohonan yang lebih rinci tidak menghapus garis
dasar sebelumnya; keduanya disimpan sebagai konfigurasi bukti yang dapat
dibandingkan.

\subsection{3.5.7 Etika, lisensi, dan
keamanan}\label{etika-lisensi-dan-keamanan}

Data terbatas, berpotensi mengandung informasi sensitif, atau memiliki
syarat penggunaan disimpan terpisah dari keluaran WebGL publik. Peta
publik hanya memuat data yang boleh dipublikasikan, agregasi yang sah,
atau proksi yang diberi label. Nilai penggeser tidak boleh mengungkap
data mentah atau memungkinkan rekonstruksi unit sensitif.

\section{3.6 Metode analisis data}\label{metode-analisis-data}

\subsection{3.6.1 Langkah 1: audit, dokumentasi, dan pembekuan garis
dasar}\label{langkah-1-audit-dokumentasi-dan-pembekuan-garis-dasar}

Analisis dimulai dengan tabel audit yang menghubungkan konstruk,
indikator, sumber, tahun, unit, polaritas, status bukti, aturan
transformasi, dan batas klaim. Garis dasar dibekukan setelah tahap audit
awal agar skenario tidak mengubah nilai observasi.

Pada tahap ini ditetapkan pula indikator yang dikeluarkan karena
duplikasi, cakupan yang tidak memadai, proporsi nilai hilang yang
mengganggu keterbandingan, ketidakcocokan unit, atau ketidakjelasan
definisi. Pengeluaran indikator dicatat sebagai keputusan metodologis,
bukan dihapus tanpa jejak.

\subsection{3.6.2 Langkah 2: delineasi pusat sekunder dan wilayah
tangkapan
fungsional}\label{langkah-2-delineasi-pusat-sekunder-dan-wilayah-tangkapan-fungsional}

Kandidat pusat sekunder berangkat dari daftar yang dikenali dalam sumber
resmi, kemudian diuji dengan bukti fungsi dan hubungan yang tersedia.
Bukti yang dapat digunakan meliputi kepadatan aktivitas, fasilitas,
pekerjaan, kawasan industri, posisi jaringan, arus perjalanan, atau
kombinasi yang ditetapkan sebelum pengukuran fungsi relatif dan
klasifikasi.

Setiap kandidat diberi alasan masuk, sumber, unit, periode, dan tingkat
ketidakpastian. Aturan pembentukan wilayah tangkapan fungsional
dibekukan untuk setiap konfigurasi bukti sebelum hasil fungsi relatif
dihitung. Wilayah tangkapan dapat berbeda antarkonfigurasi bukti, tetapi
perubahan tersebut harus diberi versi dan tidak boleh ditampilkan
sebagai perubahan kondisi metropolitan. Wilayah tangkapan merupakan
konstruksi analitis, bukan batas administratif baru.

\subsection{3.6.3 Langkah 3: pengukuran integrasi
metropolitan}\label{langkah-3-pengukuran-integrasi-metropolitan}

Integrasi diukur menurut sumbu bukti mobilitas dan hubungan metropolitan
berikut.

{\def\LTcaptype{none} % do not increment counter
\begin{longtable}[]{@{}
  >{\raggedright\arraybackslash}p{(\linewidth - 4\tabcolsep) * \real{0.3333}}
  >{\raggedright\arraybackslash}p{(\linewidth - 4\tabcolsep) * \real{0.3333}}
  >{\raggedright\arraybackslash}p{(\linewidth - 4\tabcolsep) * \real{0.3333}}@{}}
\toprule\noalign{}
\begin{minipage}[b]{\linewidth}\raggedright
Tingkat
\end{minipage} & \begin{minipage}[b]{\linewidth}\raggedright
Bukti utama
\end{minipage} & \begin{minipage}[b]{\linewidth}\raggedright
Keluaran yang sah
\end{minipage} \\
\midrule\noalign{}
\endhead
\bottomrule\noalign{}
\endlastfoot
M3 & OD antarkecamatan, zona analisis lalu lintas (TAZ), atau antarpusat
yang kompatibel & orientasi ke inti, relasi antarpusat,
\emph{inflow-outflow}, dan \emph{self-containment} bila definisi
mendukung \\
M2 & orientasi komuter, OD kabupaten/kota, atau tabulasi agregat &
orientasi dan beban pada resolusi sumber; tidak mengarang OD
antarkecamatan \\
M1 & OD sintetis, marginal bangkitan-tarikan, jaringan, waktu tempuh &
potensi interaksi dan aksesibilitas; bukan integrasi aktual \\
M0 & tidak ada relasi atau proksi yang dapat dipertanggungjawabkan &
modul integrasi dikeluarkan \\
\end{longtable}
}

Jika tersedia, arah arus dibedakan dari kedekatan. Waktu tempuh atau
jarak jaringan saja hanya mengukur aksesibilitas potensial. Persentase
komuter ke inti dapat mengukur orientasi pada unit sumber, tetapi tidak
otomatis menjelaskan hubungan antarpusat sekunder.

Jika matriks asal-tujuan yang kompatibel tersedia, proporsi orientasi
unit \(i\) menuju inti Jakarta dihitung sebagai

\[
O_i = \frac{\sum_{j \in J} F_{ij}}{\sum_j F_{ij}},
\]

dengan \(F_{ij}\) sebagai arus dari unit \(i\) menuju unit \(j\) dan
\(J\) sebagai himpunan zona inti Jakarta yang ditetapkan sebelum
perhitungan. Arus internal atau \emph{self-containment} hanya dihitung
sebagai \(F_{ii}/\sum_j F_{ij}\) apabila definisi perjalanan, zona, dan
populasi pembentuk matriks mendukung. Jika sumber hanya menerbitkan
persentase orientasi, nilai sumber digunakan pada unit publikasinya
tanpa merekonstruksi pasangan arus yang tidak tersedia.

Jika bukti yang tersedia hanya jaringan atau waktu tempuh, aksesibilitas
potensial dapat dihitung dengan bentuk umum

\[
A_i = \sum_j W_j f(c_{ij}),
\]

dengan \(W_j\) sebagai besaran kesempatan pada tujuan \(j\), \(c_{ij}\)
sebagai biaya perjalanan jaringan, dan \(f(\cdot)\) sebagai fungsi
hambatan yang dinyatakan serta diuji sensitivitasnya. Nilai \(A_i\)
tidak disebut integrasi aktual karena tidak memuat arus yang teramati.

Profil integrasi tidak harus berupa satu indeks. Bila indeks ringkas
digunakan, setiap komponen, polaritas, bobot, dan aturan normalisasi
disimpan. Integrasi yang tinggi tidak otomatis berarti manfaat tinggi
karena dapat berdampingan dengan beban dan ketergantungan.

\subsection{3.6.4 Langkah 4: pengukuran fungsi lokal dan fungsi
relatif}\label{langkah-4-pengukuran-fungsi-lokal-dan-fungsi-relatif}

Fungsi lokal diukur menurut sumbu kematangan dan fungsi berikut.

{\def\LTcaptype{none} % do not increment counter
\begin{longtable}[]{@{}
  >{\raggedright\arraybackslash}p{(\linewidth - 4\tabcolsep) * \real{0.3333}}
  >{\raggedright\arraybackslash}p{(\linewidth - 4\tabcolsep) * \real{0.3333}}
  >{\raggedright\arraybackslash}p{(\linewidth - 4\tabcolsep) * \real{0.3333}}@{}}
\toprule\noalign{}
\begin{minipage}[b]{\linewidth}\raggedright
Tingkat
\end{minipage} & \begin{minipage}[b]{\linewidth}\raggedright
Bukti utama
\end{minipage} & \begin{minipage}[b]{\linewidth}\raggedright
Keluaran yang sah
\end{minipage} \\
\midrule\noalign{}
\endhead
\bottomrule\noalign{}
\endlastfoot
U3 & layanan atau Podes kompatibel dengan pekerjaan menurut lokasi kerja
& tolok ukur layanan dan ekonomi dengan cakupan dilaporkan \\
U2 & satu dimensi langsung, misalnya layanan teramati tetapi pekerjaan
agregat & kematangan pada dimensi terukur; kematangan ekonomi penuh
tidak diklaim \\
U1 & titik minat, direktori, morfologi, kawasan industri, atau konteks
agregat & kehadiran/proksi fungsi dengan audit cakupan dan duplikasi \\
U0 & tidak ada indikator fungsi yang sebanding & modul fungsi relatif
dikeluarkan \\
\end{longtable}
}

Fungsi relatif dihitung terpisah untuk setiap konstruk \(k\), misalnya
fungsi layanan atau pekerjaan, agar residual pada satu konstruk tidak
diberi nama konstruk lain. Secara umum, tingkat fungsi yang diharapkan
dapat dinyatakan sebagai

\[
g_k(Y_{ik}) = \alpha_k + \beta_k h(S_i) + \boldsymbol{\gamma}_k^{\mathsf T}\mathbf{X}_i + \varepsilon_{ik},
\]

dengan \(Y_{ik}\) sebagai fungsi aktual pada konstruk \(k\), \(S_i\)
sebagai ukuran lokal, \(\mathbf{X}_i\) sebagai karakteristik lokal yang
ditetapkan sebelum klasifikasi, serta \(g_k(\cdot)\) dan \(h(\cdot)\)
sebagai transformasi yang dipilih berdasarkan skala pengukuran serta
diagnostik data. Residual fungsi dihitung sebagai

\[
R_{ik}=g_k(Y_{ik})-\widehat{g_k(Y_{ik})}.
\]

Residual positif berarti surplus relatif dan residual negatif berarti
defisit relatif pada konstruk \(k\). Residual layanan tidak disebut
surplus produktivitas; nilai titik minat tidak disebut pekerjaan; dan
potret satu tahun tidak disebut \emph{functional upgrading} temporal.

Pemilihan tolok ukur dilakukan sebelum tipologi dibentuk. Model
fungsi-ukuran digunakan hanya jika jumlah unit efektif, bentuk
distribusi, residual, dan pengaruh pencilan memadai untuk konstruk yang
dimodelkan. Jika syarat tersebut tidak terpenuhi, tingkat yang
diharapkan dibentuk melalui kelompok pembanding ukuran yang transparan
atau ringkasan pembanding lain yang dapat direproduksi. Pilihan model,
populasi pembanding, transformasi, prediksi, dan ketidakpastian dicatat.
Jika tolok ukur hanya memakai pusat internal kawasan studi, hasil
disebut relatif terhadap kawasan studi dan tidak digeneralisasikan
sebagai ukuran universal.

\subsection{3.6.5 Langkah 5: pemisahan manfaat, ketergantungan, dan
beban}\label{langkah-5-pemisahan-manfaat-ketergantungan-dan-beban}

Manfaat dan beban tidak dipaksa masuk ke satu sumbu yang saling
meniadakan. Profil manfaat dibentuk dari akses pekerjaan, surplus
fungsi, produktivitas atau upah, keragaman fungsi, dan kemampuan
memenuhi kebutuhan lokal hanya jika indikator tersebut benar-benar
teramati atau memiliki proksi yang diberi label. Profil beban dibentuk
dari orientasi yang sangat bergantung pada inti, durasi atau biaya
perjalanan, dan defisit fungsi pada unit yang kompatibel. Kebocoran
nilai, \emph{upgrading}, atau spesialisasi subordinat tidak menjadi
indikator operasional apabila data yang diperlukan tidak tersedia;
istilah tersebut hanya dapat digunakan sebagai interpretasi terbatas
yang ditopang bukti.

Komuter dapat berarti akses yang meningkat sekaligus beban perjalanan.
Industrialisasi dapat berarti \emph{upgrading} atau fungsi subordinat.
Oleh karena itu, polaritas ditetapkan per konstruk dan tidak disimpulkan
dari tanda positif atau negatif mentah.

\subsection{3.6.6 Langkah 6: diagnosis tipologi dan
gradien}\label{langkah-6-diagnosis-tipologi-dan-gradien}

Diagnosis dimulai dari vektor profil multidimensi, kemudian diringkas
menjadi kategori jika diperlukan untuk komunikasi. Jika beberapa
indikator digabung dalam satu dimensi \(d\), skor ringkas hanya dihitung
dalam satu konfigurasi bukti dengan himpunan indikator yang tetap:

\[
Z_{id}=\frac{\sum_{k \in K_d} w_k z_{ik}}{\sum_{k \in K_d} w_k}, \qquad w_k \geq 0.
\]

Nilai \(z_{ik}\) adalah indikator yang telah ditransformasi dan
diarahkan secara konsisten, sedangkan \(w_k\) adalah bobot yang disimpan
dalam manifest. Nilai hilang tidak diubah menjadi nol. Jika komponen
wajib tidak tersedia pada satu unit, skor dimensi dan kelas unit
tersebut dinyatakan tidak dapat dihitung dalam konfigurasi itu.
Normalisasi dilakukan terhadap populasi pembanding garis dasar yang
dibekukan agar perubahan skenario tidak mengubah titik acuan secara
tersembunyi.

Logika kategori kerja ditetapkan sebagai berikut. Nilai ambang belum
diberi angka sebelum distribusi data, dasar teoretis atau kebijakan, dan
diagnostik diperiksa; ambang garis dasar yang dipilih harus dibekukan
sebelum klasifikasi dan diuji sensitivitasnya.

{\def\LTcaptype{none} % do not increment counter
\begin{longtable}[]{@{}
  >{\raggedright\arraybackslash}p{(\linewidth - 4\tabcolsep) * \real{0.3333}}
  >{\raggedright\arraybackslash}p{(\linewidth - 4\tabcolsep) * \real{0.3333}}
  >{\raggedright\arraybackslash}p{(\linewidth - 4\tabcolsep) * \real{0.3333}}@{}}
\toprule\noalign{}
\begin{minipage}[b]{\linewidth}\raggedright
Kategori kerja
\end{minipage} & \begin{minipage}[b]{\linewidth}\raggedright
Konfigurasi minimum
\end{minipage} & \begin{minipage}[b]{\linewidth}\raggedright
Batas interpretasi
\end{minipage} \\
\midrule\noalign{}
\endhead
\bottomrule\noalign{}
\endlastfoot
Relatif independen & hubungan dengan inti berada di bawah ambang
integrasi, sedangkan fungsi lokal tidak menunjukkan defisit yang berarti
& bukan bukti autarki atau tidak adanya hubungan metropolitan \\
Pola konsisten dengan \emph{borrowed size} & bukti relasional berada di
atas ambang dan residual fungsi pada konstruk yang diuji positif &
istilah ketat hanya digunakan setelah empat gerbang bukti terpenuhi;
``integrasi produktif'' bukan sinonim otomatis \\
Campuran atau transisional & arah residual berbeda antarkonstruk, atau
manfaat dan beban sama-sama tinggi & perbedaan dimensi harus tetap
terlihat dan tidak dirata-ratakan menjadi label tunggal \\
Pola konsisten dengan \emph{agglomeration shadow} & bukti relasional
berada di atas ambang dan residual fungsi pada konstruk yang diuji
negatif; bukti beban dapat memperkuat diagnosis & tidak membuktikan
bahwa Jakarta menyebabkan defisit tersebut \\
Periferi lemah & hubungan dengan inti dan fungsi lokal sama-sama rendah
& tidak disebut \emph{shadow} Jakarta tanpa hubungan relasional yang
memadai \\
Tidak terselesaikan & unit tidak kompatibel, komponen wajib hilang, atau
kelas tidak stabil pada spesifikasi utama & ketidakpastian dilaporkan
sebagai hasil, bukan dipaksa masuk ke kelas terdekat \\
\end{longtable}
}

Label ketat \emph{borrowed size} atau \emph{agglomeration shadow} hanya
dapat dipakai bila empat gerbang berikut terpenuhi:

\begin{enumerate}
\def\labelenumi{\arabic{enumi}.}
\tightlist
\item
  ada bukti relasional terhadap Jakarta, bukan hanya jarak atau waktu
  tempuh;
\item
  ada tolok ukur fungsi relatif yang mengukur konstruk substantif pada
  unit yang memadai;
\item
  bukti relasional dan fungsi dapat dihubungkan pada unit, wilayah
  tangkapan, dan periode yang sebanding; dan
\item
  label tidak runtuh pada uji sensitivitas utama atau tidak bergantung
  pada satu proksi bermasalah.
\end{enumerate}

Jika satu gerbang gagal, istilah yang dipakai adalah \textbf{profil},
\textbf{surplus/defisit relatif}, atau \textbf{pola yang konsisten
dengan} konstruk tertentu. Penurunan klaim menjadi bagian dari hasil
metodologis.

\subsection{3.6.7 Langkah 7: diagnostik spasial dan ketahanan
hasil}\label{langkah-7-diagnostik-spasial-dan-ketahanan-hasil}

Paket diagnostik minimum mencakup korelasi dan duplikasi indikator,
pengeluaran satu indikator secara bergiliran
(\emph{leave-one-indicator-out}), variasi bobot dan ambang, serta
perubahan unit ketika perbandingan multiresolusi dapat dilakukan secara
sah. Perubahan status mutu bukti dari D ke A, S, atau P diuji sebagai
konfigurasi bukti terpisah. Perubahan definisi pusat atau wilayah
tangkapan diuji apabila lebih dari satu delineasi dapat
dipertanggungjawabkan.

Autokorelasi global dan pengelompokan lokal hanya dihitung apabila
jumlah unit, geometri, dan matriks bobot spasial memadai. Definisi
ketetanggaan atau jarak serta penanganan pulau spasial dinyatakan
sebelum perhitungan dan diuji dengan spesifikasi alternatif yang wajar.
Model regresi spasial bukan syarat desain minimum S-1; model tersebut
hanya digunakan sebagai penguatan jika diagnostik menunjukkan kebutuhan
dan asumsi model dapat dipenuhi.

Ketahanan tidak diringkas hanya sebagai persentase kelas yang sama. Peta
perpindahan, rentang skor, unit tidak terselesaikan, dan alasan
perubahan perlu ditampilkan. Kategori ``tidak terselesaikan'' adalah
keluaran yang sah ketika spesifikasi menghasilkan kesimpulan berbeda
secara substantif.

\subsection{3.6.8 Langkah 8: skenario, sensitivitas, dan implementasi
WebGL}\label{langkah-8-skenario-sensitivitas-dan-implementasi-webgl}

Analisis membedakan tiga keluarga perubahan.

\begin{enumerate}
\def\labelenumi{\arabic{enumi}.}
\tightlist
\item
  \textbf{Skenario bukti (\texttt{EV-*})} menguji hasil ketika sumber
  atau status mutu bukti berubah. Setiap skenario menyimpan sumber,
  tahun, unit, status D/A/S/P/∅, indikator yang aktif, dan klaim yang
  diperbolehkan. Perbedaan kelas antarskenario bukti menunjukkan
  ketergantungan kesimpulan pada bukti, bukan perubahan kondisi
  metropolitan.
\item
  \textbf{Skenario parameter (\texttt{P-*})} menguji bobot, ambang,
  normalisasi, indikator, atau spesifikasi lain dalam satu konfigurasi
  bukti yang sama. Skenario ini merupakan analisis sensitivitas model.
\item
  \textbf{Skenario nilai indikator} mengubah nilai dalam rentang
  terkalibrasi untuk pertanyaan \emph{what-if}. Skenario ini merupakan
  eksplorasi kondisional, bukan bagian dari diagnosis \emph{status quo}
  atau prediksi kausal kebijakan.
\end{enumerate}

Garis dasar observasi, konfigurasi bukti, dan populasi pembanding
dibekukan sebelum ketiga keluarga perubahan dijalankan. Setiap
konfigurasi menyimpan nilai garis dasar, nilai skenario, selisih, kelas
sebelum-sesudah, dan status kelayakan. Nilai yang tidak teramati tidak
diberi tampilan seolah-olah merupakan observasi, dan kombinasi yang
tidak masuk akal secara substantif ditolak oleh aturan validasi. Batas
administratif DKI selalu terlihat sebagai lapisan tetap.

Keluaran empiris utama-tabel, profil, peta diagnosis, dan audit
ketahanan-harus dapat dihasilkan tanpa WebGL. WebGL menerapkan model
yang sama sebagai antarmuka keluaran kedua. Keluaran minimalnya adalah
peta garis dasar, peta skenario, peta selisih, daftar unit yang
berpindah kelas, dan peta ketahanan atau ketidakpastian. Penghitungan
ulang serta visualisasi interaktif tidak disebut data waktu nyata atau
prediksi masa depan.

\subsection{3.6.9 Matriks keterlacakan pertanyaan, bukti, metode, dan
keluaran}\label{matriks-keterlacakan-pertanyaan-bukti-metode-dan-keluaran}

{\def\LTcaptype{none} % do not increment counter
\begin{longtable}[]{@{}
  >{\raggedright\arraybackslash}p{(\linewidth - 8\tabcolsep) * \real{0.2000}}
  >{\raggedright\arraybackslash}p{(\linewidth - 8\tabcolsep) * \real{0.2000}}
  >{\raggedright\arraybackslash}p{(\linewidth - 8\tabcolsep) * \real{0.2000}}
  >{\raggedright\arraybackslash}p{(\linewidth - 8\tabcolsep) * \real{0.2000}}
  >{\raggedright\arraybackslash}p{(\linewidth - 8\tabcolsep) * \real{0.2000}}@{}}
\toprule\noalign{}
\begin{minipage}[b]{\linewidth}\raggedright
Pertanyaan
\end{minipage} & \begin{minipage}[b]{\linewidth}\raggedright
Konstruk dan bukti minimum
\end{minipage} & \begin{minipage}[b]{\linewidth}\raggedright
Unit perhitungan
\end{minipage} & \begin{minipage}[b]{\linewidth}\raggedright
Operasi utama
\end{minipage} & \begin{minipage}[b]{\linewidth}\raggedright
Keluaran dan batas klaim
\end{minipage} \\
\midrule\noalign{}
\endhead
\bottomrule\noalign{}
\endlastfoot
Q1: hubungan metropolitan & arus atau orientasi terarah M2; M3 jika OD
rinci kompatibel; M1 hanya untuk aksesibilitas potensial & unit
asal-tujuan atau unit publikasi sumber & proporsi orientasi, profil
\emph{inflow-outflow} atau \emph{self-containment} jika sah, serta
aksesibilitas jaringan sebagai modul terpisah & peta dan profil
integrasi/orientasi; M1 hanya menghasilkan aksesibilitas potensial \\
Q2: fungsi dan kinerja relatif & fungsi lokal U2; U3 jika layanan dan
pekerjaan kompatibel; U1 hanya proksi fungsi & unit fungsi yang dapat
dipadankan dengan pusat atau wilayah tangkapan & tolok ukur
fungsi-ukuran per konstruk, residual, dan ketidakpastian & peta fungsi
aktual dan surplus/defisit relatif; residual layanan bukan
produktivitas \\
Q3: diagnosis relasional & bukti Q1 dan Q2 pada unit serta periode yang
sebanding, ditambah indikator manfaat/beban yang lolos audit & pusat
sekunder-wilayah tangkapan & profil multidimensi, gradien, logika
kategori, dan empat gerbang bukti & atlas \emph{status quo}; label ketat
hanya untuk unit yang memenuhi seluruh gerbang \\
Q4: ketahanan diagnosis & garis dasar beku, skenario bukti, dan skenario
parameter terdokumentasi & sama dengan unit keluaran garis dasar atau
agregasi yang sah & pengeluaran indikator, variasi bobot/ambang,
perubahan unit/delineasi, dan perbandingan konfigurasi bukti & rentang
skor, perpindahan kelas, peta ketahanan, dan unit tidak terselesaikan;
bukan prediksi kebijakan \\
\end{longtable}
}

Indikator final tidak dipilih hanya karena tersedia. Setiap indikator
harus memiliki definisi konstruk, sumber, tahun, unit, polaritas, status
bukti, transformasi, aturan nilai hilang, bobot jika digunakan, dan
batas klaim. Keluaran analitis serta visualisasi komunikasi memakai
model dan parameter yang sama; visualisasi tidak menghaluskan data
melebihi resolusi aslinya.

\subsection{3.6.10 Desain minimum dan
penguatan}\label{desain-minimum-dan-penguatan}

{\def\LTcaptype{none} % do not increment counter
\begin{longtable}[]{@{}
  >{\raggedright\arraybackslash}p{(\linewidth - 4\tabcolsep) * \real{0.3333}}
  >{\raggedright\arraybackslash}p{(\linewidth - 4\tabcolsep) * \real{0.3333}}
  >{\raggedright\arraybackslash}p{(\linewidth - 4\tabcolsep) * \real{0.3333}}@{}}
\toprule\noalign{}
\begin{minipage}[b]{\linewidth}\raggedright
Tingkat
\end{minipage} & \begin{minipage}[b]{\linewidth}\raggedright
Isi desain
\end{minipage} & \begin{minipage}[b]{\linewidth}\raggedright
Klaim yang diperbolehkan
\end{minipage} \\
\midrule\noalign{}
\endhead
\bottomrule\noalign{}
\endlastfoot
Minimum yang masih menjawab tesis secara terbatas & satu periode
sebanding; bukti orientasi terarah sekurang-kurangnya M2; fungsi lokal
langsung sekurang-kurangnya U2; unit yang dapat dipadankan; tolok ukur
fungsi; gradien; dan sensitivitas indikator, bobot, serta ambang & pola
deskriptif-relasional yang konsisten atau tidak konsisten dengan
\emph{borrowed size} dan \emph{agglomeration shadow} pada konstruk yang
diukur \\
Garis dasar proksi & aksesibilitas atau interaksi sintetis M1 dan fungsi
proksi U1 & profil aksesibilitas-fungsi dan pemetaan awal; tidak
menghasilkan diagnosis relasional ketat \\
Penguatan & OD rinci, pekerjaan menurut lokasi kerja, fungsi lokal
kompatibel, beberapa periode, validasi silang sumber, dan diagnostik
spasial yang memenuhi syarat & diagnosis fungsi/kinerja relatif yang
lebih kuat pada cakupan yang terlayani; klaim kausal tetap memerlukan
desain identifikasi tambahan \\
\end{longtable}
}

Jika hanya fungsi layanan yang tersedia, istilah ``kematangan''
dipersempit menjadi kematangan layanan atau kehadiran fungsi. Jika bukti
arus aktual tidak tersedia, istilah ``integrasi'' dipersempit menjadi
orientasi atau aksesibilitas sesuai sumber. Pembuatan WebGL tidak
menentukan kelayakan empiris; antarmuka hanya dibangun dari keluaran
analisis yang sah.

\subsection{3.6.11 Validitas, reliabilitas, dan
reproduktibilitas}\label{validitas-reliabilitas-dan-reproduktibilitas}

\textbf{Validitas konstruk} dijaga dengan memisahkan ukuran lokal,
hubungan relasional, aksesibilitas potensial, fungsi, kinerja, manfaat,
dan beban. Titik minat, waktu tempuh, atau jumlah fasilitas tidak
langsung diperlakukan sebagai produktivitas, integrasi aktual, atau
kematangan ekonomi.

\textbf{Validitas analitis dan reliabilitas hasil} dijaga melalui aturan
transformasi yang terdokumentasi, tolok ukur yang ditetapkan sebelum
klasifikasi, pemeriksaan duplikasi indikator, pengeluaran satu indikator
secara bergiliran, serta variasi bobot, ambang, unit, dan status bukti.
Penghitungan dengan konfigurasi masukan serta parameter yang sama harus
menghasilkan keluaran yang sama. Ketidaksesuaian antarsumber dilaporkan
dan tidak diselesaikan dengan memilih sumber yang menghasilkan
klasifikasi yang diinginkan.

\textbf{Validitas eksternal} dibatasi pada unit, tahun, dan kawasan yang
benar-benar tercakup. Hasil kabupaten/kota tidak digeneralisasikan ke
kecamatan; hasil pusat terpilih dari Disnaker atau DPMPTSP tidak
digeneralisasikan ke seluruh kawasan; dan pola dalam sistem Jakarta
tidak otomatis berlaku untuk metropolitan lain.

Paket reproduksi minimal memuat:

\begin{enumerate}
\def\labelenumi{\arabic{enumi}.}
\tightlist
\item
  kamus data, log akses, dan registri asal-usul data;
\item
  geometri serta tabel korespondensi versi analisis;
\item
  skrip atau langkah transformasi dan analisis;
\item
  nilai garis dasar dan parameter skenario;
\item
  tabel indikator, bobot, ambang, dan versi model;
\item
  daftar data yang dikeluarkan beserta alasannya; dan
\item
  manifest keluaran publik serta keluaran terbatas.
\end{enumerate}

\subsection{3.6.12 Aturan keputusan, keluaran, dan batas
klaim}\label{aturan-keputusan-keluaran-dan-batas-klaim}

Ketika jawaban permohonan data diterima, keputusan dilakukan melalui
urutan berikut:

\begin{enumerate}
\def\labelenumi{\arabic{enumi}.}
\tightlist
\item
  memeriksa isi berkas dan metadata;
\item
  menetapkan status D, A, S, P, atau ∅ berdasarkan isi, bukan nama
  produk;
\item
  memeriksa kecocokan unit, tahun, definisi, cakupan, dan reliabilitas;
\item
  memperbarui tingkat bukti M dan U;
\item
  memilih keluarga metodologis yang sesuai;
\item
  memperbarui indikator dan tolok ukur tanpa mengganti pertanyaan
  substantif;
\item
  menjalankan garis dasar terbuka dan konfigurasi penguatan secara
  paralel; serta
\item
  mengunci klaim yang dapat dibuat sebelum menulis interpretasi.
\end{enumerate}

Jumlah pintu data yang terbuka bukan ukuran mutu metode. Satu berkas
berstatus D pada konstruk yang tepat dapat lebih berguna daripada
beberapa berkas agregat yang tidak kompatibel. Perubahan berkas atau
status menghasilkan versi konfigurasi bukti baru tanpa menimpa manifest
sebelumnya.

Keluaran utama penelitian adalah tabel audit dan manifest mutu bukti;
peta serta profil hubungan metropolitan atau aksesibilitas; peta fungsi
aktual dan fungsi relatif; profil manfaat serta beban; gradien atau
tipologi \emph{status quo}; tabel dan peta sensitivitas; serta
penjelajah WebGL garis dasar-skenario-selisih-ketahanan. Hasil juga
dilaporkan menurut lapisan administrasi untuk membantu interpretasi
tanggung jawab tanpa mengubah batas kewenangan.

Dengan desain minimum, penelitian dapat menyatakan pola hubungan atau
aksesibilitas, fungsi lokal yang teramati, surplus atau defisit relatif
pada konstruk yang diukur, tipologi deskriptif, gradien, dan stabilitas
hasil terhadap asumsi yang diuji. Penelitian tidak dapat menyatakan
bahwa Jakarta secara kausal menyebabkan ketergantungan, bahwa satu
radius merupakan batas metropolitan yang benar, bahwa pusat sekunder
pasti mengalami \emph{upgrading}, atau bahwa perubahan nilai dalam WebGL
memprediksi dampak kebijakan masa depan.

Klaim \emph{borrowed size} dan \emph{agglomeration shadow} digunakan
secara ketat hanya setelah empat gerbang bukti terpenuhi. Jika tidak,
hasil ditulis sebagai pola yang ``konsisten dengan'' atau ``tidak
konsisten dengan'' konstruk tersebut pada unit, periode, dan konfigurasi
bukti yang disebutkan.

\end{document}
